problem:  
Let \((M^n,g,f)\) be a complete smooth metric measure space with measure \(d\mu = e^{-f}d\mathrm{vol}_g\) and suppose \((M,d,\mu)\) satisfies volume doubling and a scale-invariant Poincaré inequality.  
Assume \(M\) has one nonparabolic end \(E\subset M\) that is **asymptotically conical in the weighted sense**, meaning there exists a compact set \(K\) and a diffeomorphism  
\[
E\setminus K \cong (R_0,\infty)\times S^{n-1}
\]
with metric  
\[
g = dr^2 + r^2 h_r, \qquad \|h_r - h_\infty\|_{C^2(S^{n-1})} = O(r^{-\alpha}),
\]
and weight \(f=f(r,\theta)\) satisfying
\[
|f(r,\theta) - \phi(r)| + r |\partial_r(f-\phi)| \le C\, r^{-\alpha}, \quad \text{for some radial } \phi(r).
\]

Let \(\mathcal{H}_f^\infty(M)\) be the space of bounded \(f\)-harmonic functions, i.e.
\[
\mathcal{H}_f^\infty(M) = \{u \in L^\infty(M,\mu)\mid \Delta_f u = 0\}, \quad \Delta_f = \Delta - \langle\nabla f,\nabla\cdot\rangle.
\]

---

**Open problem:**  
Does the Liouville property
\[
\mathcal{H}_f^\infty(M) = \mathbb{R}
\]
remain stable under all \(O(r^{-\alpha})\)-asymptotically conical perturbations of the link metric \(h_\infty\) and the drift potential \(f\), provided the Bakry–Émery tensor satisfies
\[
\mathrm{Ric}_f \ge -K r^{-2},
\]
or can one find examples with the **same asymptotic cone and weight profile** but nonconstant bounded \(f\)-harmonic functions?

Equivalently: *Is there a universal decay exponent \(\alpha_c>0\) such that Liouville-type vanishing for bounded \(f\)-harmonic functions holds for all \(O(r^{-\alpha})\)-convergent weighted cones if and only if \(\alpha>\alpha_c\)?*  
If not, can one characterize the minimal geometric obstructions at subcritical decay rates that allow new bounded harmonic modes to appear?

---

Why is it a "good" problem:  
This problem precisely isolates the frontier of **analytic and geometric stability** of Liouville-type theorems in the weighted asymptotically conical (AC) setting—a context central to analysis on manifolds, Ricci solitons, and scattering geometry.  
Known vanishing results in asymptotically conical geometry rely heavily on \(\alpha>1\) to invoke Fredholm and decay theory. Whether the bounded \(f\)-harmonic Liouville property survives at or below this critical threshold (\(\alpha\le1\)) remains completely open, even in the unweighted case.  

The problem’s simplicity (a decay rate question for geometric convergence) hides a deep analytic challenge: it demands understanding how slow metric and weight convergence can destabilize the spectral and potential-theoretic rigidity of the Laplacian. This requires tools from microlocal analysis, diffusion theory, and potential theory on manifolds with ends.  

By bridging geometric asymptotics with analytic stability of harmonic function spaces, the question lies at the conceptual heart of geometric analysis and spectral theory—the quintessential “narrow entrance, profound depth’’ problem.